\documentclass[11pt]{article}

\usepackage[a4paper,margin=2.4cm]{geometry}
\usepackage{amsmath,amssymb,mathtools}
\usepackage{algorithm}
\usepackage{algpseudocode}
\usepackage{microtype}

\newcommand{\R}{\mathbb{R}}
\newcommand{\col}{\operatorname{col}}
\newcommand{\blkdiag}{\operatorname{blkdiag}}
\newcommand{\argmin}{\operatorname*{arg\,min}}

\title{A Structured Dual Active-Set Solver for Linear--Quadratic MPC}
\author{}
\date{}

\begin{document}
\maketitle

\begin{abstract}
This note describes a dual active-set solver for linear--quadratic
optimal-control problems with linear state and input constraints.
The dual Hessian admits a Gram factorization that can be constructed
from the Riccati recursion, so neither the condensed Hessian nor its
inverse needs to be formed explicitly. The Gram factor is used to
maintain the active dual Hessian through recursive \(LDL^\top\)
updates. The complete vector of primal slacks, however, is evaluated
by solving an unconstrained LQR problem whose linear terms depend on
the current dual multipliers. For fixed stage dimensions, this slack
evaluation has complexity linear in the prediction horizon.
\end{abstract}

\section{Optimal-control problem and condensation}

Consider
\begin{equation}
\begin{aligned}
\min_{\{x_t,u_t\}}\quad&
 \frac{1}{2}x_N^\top P_Nx_N+p_N^\top x_N\\
&+\sum_{t=0}^{N-1}
 \left(
 \frac{1}{2}x_t^\top Qx_t+q_t^\top x_t
 +\frac{1}{2}u_t^\top Ru_t+r_t^\top u_t
 \right)\\
\text{s.t.}\quad&
 x_{t+1}=Ax_t+Bu_t,\qquad t=0,\ldots,N-1,\\
&Cu_t\leq c,\qquad t=0,\ldots,N-1,\\
&Dx_t\leq d,\qquad t=1,\ldots,N,\\
&x_0\ \text{given},
\end{aligned}
\label{eq:ocp}
\end{equation}
where
\[
R=R^\top\succ0,\qquad
Q=Q^\top\succeq0,\qquad
P_N=P_N^\top\succeq0.
\]
The state and input dimensions are denoted by \(n_x\) and \(n_u\),
while \(m_x\) and \(m_u\) denote the numbers of state and input
inequalities at each stage.

Define
\[
\bm{x}=\col(x_1,\ldots,x_N),\qquad
\bm{u}=\col(u_0,\ldots,u_{N-1}).
\]
Eliminating the dynamics gives
\begin{equation}
\bm{x}=\mathcal A x_0+\mathcal B\bm{u},
\label{eq:state-condensing}
\end{equation}
where, using zero-based block indices,
\[
[\mathcal A]_k=A^{k+1},\qquad
[\mathcal B]_{k,t}=
\begin{cases}
A^{k-t}B,&t\leq k,\\
0,&t>k.
\end{cases}
\]

Introduce
\[
\begin{aligned}
\bar Q&=\blkdiag(
 \underbrace{Q,\ldots,Q}_{N-1\text{ blocks}},P_N),\\
\bar R&=I_N\otimes R,\\
\bar q&=\col(q_1,\ldots,q_{N-1},p_N),\\
\bar r&=\col(r_0,\ldots,r_{N-1}),\\
\bar C&=I_N\otimes C,\qquad
\bar D=I_N\otimes D,\\
\bar c&=\bm{1}_N\otimes c,\qquad
\bar d=\bm{1}_N\otimes d.
\end{aligned}
\]
Time-varying constraint matrices or bounds can instead be stacked
stage by stage.

After removing terms that depend only on \(x_0\), problem
\eqref{eq:ocp} becomes
\begin{equation}
\begin{aligned}
\min_{\bm{u}}\quad&
 \frac{1}{2}\bm{u}^\top H\bm{u}+g^\top\bm{u}\\
\text{s.t.}\quad&G\bm{u}\leq b,
\end{aligned}
\label{eq:condensed-qp}
\end{equation}
where
\begin{equation}
\begin{aligned}
H&=\bar R+\mathcal B^\top\bar Q\mathcal B,\\
g&=\mathcal B^\top(\bar Q\mathcal A x_0+\bar q)+\bar r,\\
G&=
\begin{bmatrix}
\bar C\\
\bar D\mathcal B
\end{bmatrix},
\qquad
b=
\begin{bmatrix}
\bar c\\
\bar d-\bar D\mathcal A x_0
\end{bmatrix}.
\end{aligned}
\label{eq:condensed-data}
\end{equation}
Since \(R\succ0\), the condensed Hessian satisfies \(H\succ0\).

\section{Dual problem and primal slacks}

Let
\[
\nu=\col(\mu,\lambda)\geq0,
\]
where
\[
\mu=\col(\mu_0,\ldots,\mu_{N-1}),\qquad
\lambda=\col(\lambda_1,\ldots,\lambda_N)
\]
are the input- and state-constraint multipliers.

Minimizing the condensed Lagrangian over \(\bm{u}\) gives the
equivalent dual minimization problem
\begin{equation}
\min_{\nu\geq0}
\quad
\phi(\nu)
=
\frac{1}{2}\nu^\top S\nu+h^\top\nu,
\label{eq:dual}
\end{equation}
where
\begin{equation}
S=GH^{-1}G^\top,\qquad
h=b+GH^{-1}g.
\label{eq:dual-data}
\end{equation}
An additive constant has been omitted.

For a fixed multiplier vector \(\nu\), let
\begin{equation}
\bm{u}(\nu)
=
-H^{-1}(g+G^\top\nu)
\label{eq:dual-primal-map}
\end{equation}
and let \(\bm{x}(\nu)\) be the corresponding state trajectory.
Then
\begin{equation}
\nabla\phi(\nu)
=
S\nu+h
=
b-G\bm{u}(\nu).
\label{eq:gradient-slacks}
\end{equation}
Thus, the gradient of the dual objective is exactly the vector of
primal slacks:
\begin{equation}
s(\nu)
=
\col\left(
\{c-Cu_t(\nu)\}_{t=0}^{N-1},
\{d-Dx_t(\nu)\}_{t=1}^{N}
\right).
\label{eq:slack-definition}
\end{equation}

In particular, let
\((\widehat{\bm{x}},\widehat{\bm{u}})\) be the unconstrained LQR
solution, corresponding to \(\nu=0\). Then
\begin{equation}
h=s(0)=
\begin{bmatrix}
\bar c-\bar C\widehat{\bm{u}}\\
\bar d-\bar D\widehat{\bm{x}}
\end{bmatrix}.
\label{eq:dual-linear-slack}
\end{equation}
Hence the online dual linear term is obtained from one unconstrained
LQR solve.

\section{Riccati factorization of the dual Hessian}

\subsection{Quadratic Riccati recursion}

Starting from \(P_N\), compute for \(t=N-1,\ldots,0\)
\begin{equation}
\begin{aligned}
\Sigma_t&=R+B^\top P_{t+1}B=L_tL_t^\top,\\
K_t&=\Sigma_t^{-1}B^\top P_{t+1}A,\\
A_t^{\mathrm c}&=A-BK_t,\\
P_t&=Q+K_t^\top RK_t
 +(A_t^{\mathrm c})^\top P_{t+1}A_t^{\mathrm c},
\end{aligned}
\label{eq:riccati}
\end{equation}
where \(L_t\) is lower triangular. These quantities depend only on
the quadratic cost and the system dynamics and can therefore be
computed offline.

Define the closed-loop transition matrix
\[
\Phi_{k,j}=
\begin{cases}
A_{k-1}^{\mathrm c}\cdots A_j^{\mathrm c},&k>j,\\
I,&k=j.
\end{cases}
\]

\subsection{Factors of \(H^{-1}\)}

Let \(U\in\R^{Nn_u\times Nn_u}\) and
\(X\in\R^{Nn_x\times Nn_u}\) have block entries
\begin{equation}
U_{k,t}=
\begin{cases}
0,&t>k,\\
L_k^{-\top},&t=k,\\
-K_k\Phi_{k,t+1}B L_t^{-\top},&t<k,
\end{cases}
\label{eq:U-blocks}
\end{equation}
and
\begin{equation}
X_{k,t}=
\begin{cases}
\Phi_{k+1,t+1}B L_t^{-\top},&t\leq k,\\
0,&t>k.
\end{cases}
\label{eq:X-blocks}
\end{equation}
Row block \(k\) of \(X\) corresponds to state \(x_{k+1}\).

Completing squares in the Riccati recursion gives
\begin{equation}
\begin{aligned}
H^{-1}&=UU^\top,\\
\mathcal B H^{-1}\mathcal B^\top&=XX^\top,\\
H^{-1}\mathcal B^\top&=UX^\top,
\end{aligned}
\label{eq:inverse-factorizations}
\end{equation}
with \(X=\mathcal B U\). Consequently,
\begin{equation}
S=MM^\top,\qquad
M=
\begin{bmatrix}
\bar C U\\
\bar D X
\end{bmatrix}
=GU.
\label{eq:dual-gram-factor}
\end{equation}
The entries of the dual Hessian are therefore inner products between
rows of \(M\).

\subsection{Direct construction of rows of \(M\)}

Let
\[
C_{i,:}=(a_i^u)^\top,\qquad
D_{j,:}=(a_j^x)^\top.
\]
For input constraint \(i\) at time \(k\), define
\(m_{k,i}^u\in\R^{Nn_u}\), partitioned into \(n_u\)-dimensional
blocks. Set
\begin{equation}
[m_{k,i}^u]^{(t)}=0,\quad t>k,
\qquad
[m_{k,i}^u]^{(k)}=L_k^{-1}a_i^u.
\end{equation}
Initialize
\[
\xi=-K_k^\top a_i^u
\]
and, for \(t=k-1,\ldots,0\), compute
\begin{equation}
[m_{k,i}^u]^{(t)}=L_t^{-1}B^\top\xi,
\qquad
\xi\leftarrow(A_t^{\mathrm c})^\top\xi.
\label{eq:input-feature-recursion}
\end{equation}
The corresponding row of \(M\) is \((m_{k,i}^u)^\top\).

For state constraint \(j\) on \(x_{k+1}\), initialize
\[
\xi=a_j^x
\]
and, for \(t=k,\ldots,0\), compute
\begin{equation}
[m_{k,j}^x]^{(t)}=L_t^{-1}B^\top\xi,
\qquad
\xi\leftarrow(A_t^{\mathrm c})^\top\xi,
\label{eq:state-feature-recursion}
\end{equation}
with \([m_{k,j}^x]^{(t)}=0\) for \(t>k\).

For any pair of inequalities \(\ell\) and \(j\),
\begin{equation}
S_{\ell j}=m_\ell^\top m_j.
\label{eq:hessian-inner-product}
\end{equation}
These feature vectors are used to build and update the active dual
Hessian. They need not be used to evaluate the full dual gradient.

\section{An LQR oracle for the dual gradient}

The condensed expression
\[
s(\nu)=MM^\top\nu+h
\]
is mathematically correct, but evaluating it by multiplying by the
full matrix \(M\) can be unnecessarily expensive when \(N\) is large.
Instead, \(s(\nu)\) can be recovered by solving
\begin{equation}
\min_{\{x_t,u_t\}}
\mathcal L(x,u,\nu),
\label{eq:lagrangian-min}
\end{equation}
where the dynamics are enforced and
\begin{equation}
\begin{aligned}
\mathcal L(x,u,\nu)
={}&J(x,u)
+\sum_{t=0}^{N-1}\mu_t^\top(Cu_t-c)\\
&+\sum_{t=1}^{N}\lambda_t^\top(Dx_t-d).
\end{aligned}
\label{eq:ocp-lagrangian}
\end{equation}
For fixed \(\nu\), terms involving only \(c,d,\mu,\lambda\) do not
affect the minimizing trajectory. Therefore
\eqref{eq:lagrangian-min} is an unconstrained LQR problem with the
same quadratic terms as the original problem and modified linear
terms
\begin{equation}
\begin{aligned}
\widetilde r_t(\nu)
&=r_t+C^\top\mu_t,
&&t=0,\ldots,N-1,\\
\widetilde q_0(\nu)
&=q_0,\\
\widetilde q_t(\nu)
&=q_t+D^\top\lambda_t,
&&t=1,\ldots,N-1,\\
\widetilde p_N(\nu)
&=p_N+D^\top\lambda_N.
\end{aligned}
\label{eq:modified-linear-terms}
\end{equation}

Since the quadratic terms are unchanged, the precomputed matrices
\(P_t,K_t,L_t\), and \(A_t^{\mathrm c}\) can be reused. Set
\begin{equation}
\pi_N=\widetilde p_N(\nu)
\end{equation}
and perform, for \(t=N-1,\ldots,0\), the backward recursion
\begin{equation}
\begin{aligned}
\omega_t
&=L_t^{-1}\left(
 \widetilde r_t(\nu)+B^\top\pi_{t+1}
 \right),\\
v_t
&=-L_t^{-\top}\omega_t,\\
\pi_t
&=\widetilde q_t(\nu)
 +(A_t^{\mathrm c})^\top\pi_{t+1}
 -K_t^\top\widetilde r_t(\nu).
\end{aligned}
\label{eq:dual-lqr-backward}
\end{equation}
The minimizing trajectory is obtained from
\begin{equation}
\begin{aligned}
x_0(\nu)&=x_0,\\
u_t(\nu)&=-K_tx_t(\nu)+v_t,\\
x_{t+1}(\nu)&=Ax_t(\nu)+Bu_t(\nu),
\end{aligned}
\qquad t=0,\ldots,N-1.
\label{eq:dual-lqr-forward}
\end{equation}
Finally, compute
\begin{equation}
\begin{aligned}
s_t^u(\nu)&=c-Cu_t(\nu),
&&t=0,\ldots,N-1,\\
s_t^x(\nu)&=d-Dx_t(\nu),
&&t=1,\ldots,N.
\end{aligned}
\label{eq:lqr-slacks}
\end{equation}
Stacking these quantities gives \(s(\nu)=\nabla\phi(\nu)\).

For fixed \(n_x,n_u,m_x,m_u\), the backward and forward recursions
and the evaluation of \eqref{eq:lqr-slacks} have complexity linear in
\(N\). This can be substantially cheaper than forming
\(w=M^\top\nu\) and evaluating \(Mw+h\), whose cost generally grows
quadratically with the horizon because the rows of \(M\) extend over
multiple stages.

The LQR oracle also directly returns the primal trajectory associated
with the dual candidate. At the final dual solution, this trajectory
is the optimal primal solution, so no additional primal-recovery
operation is needed.

\section{Dual active-set method}

Let \(\mathcal W\) be the working set. Multipliers indexed by
\(\mathcal W\) are free, while
\[
\nu_i=0,\qquad i\notin\mathcal W.
\]
Write \(M_{\mathcal W}\) and \(h_{\mathcal W}\) for the corresponding
rows of \(M\) and entries of \(h\). The equality-constrained dual
subproblem is
\begin{equation}
\min_{\nu_{\mathcal W}}
\frac{1}{2}\nu_{\mathcal W}^\top
M_{\mathcal W}M_{\mathcal W}^\top\nu_{\mathcal W}
+h_{\mathcal W}^\top\nu_{\mathcal W}.
\label{eq:working-subproblem}
\end{equation}
Define
\[
K_{\mathcal W}:=
M_{\mathcal W}M_{\mathcal W}^\top.
\]
When \(K_{\mathcal W}\) is nonsingular, the unconstrained minimizer
of the working-set subproblem satisfies
\begin{equation}
K_{\mathcal W}\nu_{\mathcal W}^\star=-h_{\mathcal W}.
\label{eq:working-solve}
\end{equation}

\begin{algorithm}[t]
\caption{Structured dual active-set method}
\label{alg:das}
\begin{algorithmic}[1]
\Require Riccati factors, rows of \(M\) or a row-product oracle,
\(h\), a working set \(\mathcal W\), and a dual-feasible \(\nu\)
satisfying \(\nu_{\mathcal W^c}=0\)
\Loop
    \If{\(K_{\mathcal W}\) is nonsingular}
        \State Solve
        \(K_{\mathcal W}\nu_{\mathcal W}^\star=-h_{\mathcal W}\)
        using its \(LDL^\top\) factorization
        \State Set
        \(\nu_{\mathcal W^c}^\star=0\)
        \If{\(\nu_{\mathcal W}^\star\geq0\)}
            \State \(\nu\gets\nu^\star\)
            \State Solve the LQR problem
            \(\min_{x,u}\mathcal L(x,u,\nu)\)
            using \eqref{eq:dual-lqr-backward}--%
            \eqref{eq:dual-lqr-forward}
            \State Compute all slacks \(s(\nu)\) from
            \eqref{eq:lqr-slacks}
            \If{\(s_{\mathcal W^c}\geq0\)}
                \State \Return \(\nu,\mathcal W,x(\nu),u(\nu)\)
            \Else
                \State
                \(j\gets\argmin_{i\notin\mathcal W}s_i(\nu)\)
                \State
                \(\mathcal W\gets\mathcal W\cup\{j\}\)
                and update the \(LDL^\top\) factors
            \EndIf
        \Else
            \State
            \(p_{\mathcal W}\gets
            \nu_{\mathcal W}^\star-\nu_{\mathcal W}\),
            \(p_{\mathcal W^c}\gets0\)
            \State
            \(j\gets
            \argmin_{\substack{i\in\mathcal W\\p_i<0}}
            (-\nu_i/p_i)\)
            \State
            \(\alpha\gets-\nu_j/p_j\),
            \(\nu\gets\nu+\alpha p\)
            \State
            \(\mathcal W\gets\mathcal W\setminus\{j\}\)
            and downdate the \(LDL^\top\) factors
        \EndIf
    \Else
        \State Compute \(p_{\mathcal W}\neq0\) satisfying
        \(K_{\mathcal W}p_{\mathcal W}=0\)
        \State Select its sign so that
        \(h_{\mathcal W}^\top p_{\mathcal W}<0\);
        set \(p_{\mathcal W^c}=0\)
        \If{\(p_{\mathcal W}\geq0\)}
            \State \Return primal infeasibility
        \Else
            \State
            \(j\gets
            \argmin_{\substack{i\in\mathcal W\\p_i<0}}
            (-\nu_i/p_i)\)
            \State
            \(\alpha\gets-\nu_j/p_j\),
            \(\nu\gets\nu+\alpha p\)
            \State
            \(\mathcal W\gets\mathcal W\setminus\{j\}\)
            and downdate the \(LDL^\top\) factors
        \EndIf
    \EndIf
\EndLoop
\end{algorithmic}
\end{algorithm}

At a nonsingular working-set minimizer,
\[
[s(\nu^\star)]_{\mathcal W}=0
\]
in exact arithmetic. Hence only the inactive slacks must be tested.
A negative inactive slack identifies a violated primal inequality,
and the corresponding multiplier is added to the working set.

Importantly, the LQR slack oracle is invoked only when the
working-set minimizer is dual feasible. If
\(\nu_{\mathcal W}^\star\) contains a negative component, the
algorithm immediately performs the ratio test and removes a blocking
constraint without computing the full slack vector.

If \(K_{\mathcal W}\) is singular, a direction satisfying
\[
K_{\mathcal W}p_{\mathcal W}=0,\qquad
h_{\mathcal W}^\top p_{\mathcal W}<0
\]
is a descent direction. If it is also nonnegative, the dual objective
is unbounded below along the ray, certifying infeasibility of the
primal inequalities.

\section{Recursive \(LDL^\top\) factorization}

Maintain
\begin{equation}
K_{\mathcal W}
=
\mathscr L\Delta\mathscr L^\top,
\label{eq:ldl}
\end{equation}
where \(\mathscr L\) is unit lower triangular and \(\Delta\) is
diagonal and nonnegative.

\subsection{Adding a constraint}

Suppose row \(m_j^\top\) is appended to \(M_{\mathcal W}\). Compute
\[
q=M_{\mathcal W}m_j,\qquad
\gamma=m_j^\top m_j.
\]
Solve
\begin{equation}
\mathscr L\Delta\ell=q,
\qquad
\delta=\gamma-\ell^\top\Delta\ell.
\label{eq:ldl-addition}
\end{equation}
The updated factors are
\begin{equation}
\mathscr L^+
=
\begin{bmatrix}
\mathscr L&0\\
\ell^\top&1
\end{bmatrix},
\qquad
\Delta^+
=
\begin{bmatrix}
\Delta&0\\
0&\delta
\end{bmatrix}.
\label{eq:ldl-addition-factors}
\end{equation}
A zero or numerically small value of \(\delta\) indicates that the
new row is linearly dependent on the current active rows.

\subsection{Removing a constraint}

Suppose the removed constraint is at position \(i\), and partition
the factors as
\[
\mathscr L=
\begin{bmatrix}
\mathscr L_1&0&0\\
*&1&0\\
A&\ell_i&\mathscr L_2
\end{bmatrix},
\qquad
\Delta=
\begin{bmatrix}
\Delta_1&0&0\\
0&\delta_i&0\\
0&0&\Delta_2
\end{bmatrix}.
\]
After removing index \(i\), the new factors are
\[
\mathscr L^-=
\begin{bmatrix}
\mathscr L_1&0\\
A&\widetilde{\mathscr L}_2
\end{bmatrix},
\qquad
\Delta^-=
\begin{bmatrix}
\Delta_1&0\\
0&\widetilde{\Delta}_2
\end{bmatrix},
\]
where
\begin{equation}
\widetilde{\mathscr L}_2
\widetilde{\Delta}_2
\widetilde{\mathscr L}_2^\top
=
\mathscr L_2\Delta_2\mathscr L_2^\top
+\delta_i\ell_i\ell_i^\top.
\label{eq:ldl-removal}
\end{equation}
Thus, deleting a working-set constraint reduces to a rank-one
\(LDL^\top\) update of the trailing factor.

\subsection{Null direction from a zero pivot}

If \(\Delta_{ii}=0\), partition
\[
\mathscr L=
\begin{bmatrix}
\mathscr L_1&0&0\\
\ell_i^\top&1&0\\
*&*&*
\end{bmatrix}.
\]
Solve
\[
\mathscr L_1^\top\widetilde p=-\ell_i
\]
and set
\[
p_{\mathcal W}
=
\col(\widetilde p,1,0).
\]
Then
\[
K_{\mathcal W}p_{\mathcal W}=0.
\]
Its sign is selected so that
\(h_{\mathcal W}^\top p_{\mathcal W}<0\).

\section{Offline--online implementation}

The implementation uses the Gram representation and the LQR
representation for different purposes.

\paragraph{Offline phase.}
\begin{enumerate}
\item Compute \(P_t,K_t,L_t\), and \(A_t^{\mathrm c}\) through
      \eqref{eq:riccati}.
\item Generate the feature vectors \(m_{k,i}^u\) and \(m_{k,j}^x\),
      or store sufficient data to evaluate their inner products.
\item Allocate the working-set and recursive \(LDL^\top\) workspace.
\end{enumerate}

\paragraph{Online initialization.}
\begin{enumerate}
\item Set \(\nu=0\) and solve the original unconstrained LQR problem.
\item Compute
      \[
      h=s(0)
      \]
      directly from the input and state slacks.
\item Shift the previous MPC working set and multipliers to obtain a
      warm start.
\item Reuse or reconstruct the \(LDL^\top\) factors associated with
      the shifted working set.
\end{enumerate}

\paragraph{Active-set iterations.}
\begin{enumerate}
\item Use row inner products of \(M\) only to maintain
      \(K_{\mathcal W}\) and solve the working-set subproblem.
\item Whenever a dual-feasible working-set minimizer is obtained,
      evaluate all primal slacks through one LQR backward--forward
      recursion.
\item At termination, return the state and input trajectory already
      computed by the final LQR slack evaluation.
\end{enumerate}

This separation is useful for long horizons. The active-set linear
algebra depends mainly on the number of active constraints, while the
full feasibility check is performed by an \(O(N)\) optimal-control
recursion rather than by multiplication with the complete condensed
dual Hessian.

\section{Numerical considerations}

In finite precision, use separate tolerances for dual feasibility,
primal feasibility, and singularity detection:
\[
\nu_i\geq-\varepsilon_\nu,\qquad
s_i\geq-\varepsilon_{\mathrm p},\qquad
\delta\leq
\varepsilon_{\mathrm{piv}}\max\{1,\gamma\}.
\]
Since the active-set equations imply \(s_{\mathcal W}=0\), only
\(s_{\mathcal W^c}\) should normally be used to select a constraint
to add. Large residuals on \(\mathcal W\) indicate inaccurate
working-set solves or deteriorated matrix factors.

Constraint rows should be scaled to comparable norms before
constructing the feature vectors. Deterministic tie-breaking should
be used in both the most-violated-constraint test and the ratio test.
For poorly conditioned problems, the \(LDL^\top\) factors should
occasionally be reconstructed from the current working set rather
than recursively updated indefinitely.

An outer proximal-point iteration can be included by regularizing the
quadratic cost. For a fixed regularization parameter, the associated
Riccati factors and dual-Hessian feature vectors remain reusable
across the outer iterations.

Lower bounds may either be converted to upper inequalities or handled
through signed multipliers. The latter avoids duplicating some of the
slack computations for paired upper and lower bounds.

\begin{thebibliography}{1}

\bibitem{arnstrom2021}
D.~Arnstr{\"o}m, A.~Bemporad, and D.~Axehill,
``A dual active-set solver for embedded quadratic programming using
recursive \(LDL^\top\) updates,''
arXiv:2103.16236, 2021.

\end{thebibliography}

\end{document}